\documentclass[11pt]{article}
\usepackage[a4paper,margin=25mm]{geometry}
\usepackage{amsmath,amssymb,amsthm,mathtools,hyperref,enumitem,longtable}
\hypersetup{colorlinks=true,urlcolor=blue,linkcolor=blue}
\newtheorem{theorem}{Theorem}[section]
\newtheorem{lemma}[theorem]{Lemma}
\newtheorem{proposition}[theorem]{Proposition}
\newtheorem{definition}[theorem]{Definition}
\newtheorem{warning}[theorem]{Warning}
\newcommand{\ipra}{\mathsf{iPRA}}
\newcommand{\bfam}{\mathcal B}
\newcommand{\one}{\mathbf 1}
\title{Bounded Reflection and Guarded Witness Compilation in Intuitionistic Primitive Recursive Arithmetic\\\large A reusable object-theory framework and a Lean integration plan}
\author{Research note, 16 November 2026}
\date{}
\begin{document}\maketitle
\begin{abstract}
We isolate a reusable, strictly finitary compiler for bounded arithmetic formulas in intuitionistic primitive recursive arithmetic (iPRA). Every fixed bounded formula with primitive-recursive term bounds is assigned a primitive-recursive Boolean evaluator; iPRA proves its correctness intuitionistically. Bounded existential proof obligations obtain canonical least witnesses, including explicit failure flags, while bounded universal obligations compile into a finite conjunction bit. Parameter-dependent finite coverage and guarded witness composition are proved without countable choice, excluded middle for arbitrary formulas, or any appeal to actual infinity. We distinguish semantic correctness in Lean's metatheory from syntactic derivations in the embedded object calculus, give an implementation-oriented Lean plan, and identify sharp limitations of the method. This is a new integration proposal, not a claim of a kernel-checked implementation.
\end{abstract}
\tableofcontents
\section{Repository baseline and scope}
The repository \url{https://github.com/VladimirReshetnikov/ProveIt/tree/main/Logic/PrimitiveRecursiveArithmetic} defines iPRA as intuitionistic first-order arithmetic with a symbol and defining equations for every primitive-recursive function and quantifier-free induction. Its \texttt{ConciseIPRA} uses terms \texttt{PTm}, formulas \texttt{Fm} and explicit derivations \texttt{Prf PRAAx}. The separately checked \texttt{IPRADerivations} library includes order, characteristic equality, bounded search, bounded sums, pairing and finite recursion; \texttt{Bridge} supplies closed-sentence translations with Foundation arithmetic. The unification plan explicitly lists open-sequent transport, formula round-trips, and symbol comparisons as unfinished tasks. The proof-packet README describes many uncompiled frameworks and warns that their stated main claims must not be confused with checked object-theory proofs. We therefore reuse the verified arithmetic foundations rather than representing another packet as a completed formalization.

Our narrower contribution is a \emph{boundary-aware Booleanization theorem}: a precise compilation principle that is stable under bounded quantifier alternation, whose proof produces actual iPRA derivations, and whose existential branch returns a primitive-recursive witness plus a success bit. The primitive-recursive program extracted from the compiler is intended as a reusable interface for algebraic search, finite geometry, effective continuity at a specified precision, and finite proof graphs. The method does \emph{not} claim to settle open cases involving genuinely unbounded witness search or uniform bounds not known to be primitive recursive.

\section{Syntax and proof-theoretic discipline}
Fix a finite tuple of free number variables $\vec x$. Let $\bfam$ be the least class of formulas generated by atomic equalities $s(\vec x)=t(\vec x)$, where $s,t$ are primitive-recursive terms, using $\wedge,\vee,\neg$, and bounded quantifiers
\[\exists y<B(\vec x)\, A(y,\vec x),\qquad \forall y<B(\vec x)\, A(y,\vec x).\]
Here $B$ is a specified primitive-recursive term independent of the bound variable. We use $A\to C$ as an abbreviation $\neg A\vee C$ only for formulas already known decidable, or alternatively add implication as a Boolean connective and prove the corresponding restricted reflection. An arbitrary implication is \emph{not} replaced with classical implication.

Crucial distinction: quantifier-free induction in $\ipra$ applies to formulas with no object-language quantifier. For each bounded formula we shall construct a number-valued characteristic term $e_A(\vec x)$, and induction will apply to equations such as $e_A=1$. We never take induction on the original bounded-quantifier formula as a primitive axiom. Every parameter and every primitive-recursive definition is fixed syntactically before the proof.

\begin{definition}[Boolean convention]
A Boolean term $b$ satisfies $b\le 1$, with $b=1$ meaning true and $b=0$ meaning false. Put
\[\operatorname{not}(b)=1\mathbin{\dot-} b,\quad \operatorname{and}(b,c)=bc,\quad \operatorname{or}(b,c)=1-(1-b)(1-c),\]
where subtraction denotes truncated subtraction. All displayed equations involving $1-b$ are used only when $b\le1$. Define $\operatorname{eqb}(s,t)$ by a standard primitive-recursive equality test, with iPRA derivations of $\operatorname{eqb}(s,t)=1\leftrightarrow s=t$ and its complementary case.
\end{definition}

\section{Finite aggregators and their object-theory laws}
For a Boolean-valued primitive-recursive $f(i,\vec x)$, define the finite conjunction and disjunction by primitive recursion:
\begin{align*}
U_f(0,\vec x)&=1,&U_f(n+1,\vec x)&=U_f(n,\vec x) f(n,\vec x),\\
E_f(0,\vec x)&=0,&E_f(n+1,\vec x)&=\operatorname{or}(E_f(n,\vec x),f(n,\vec x)).
\end{align*}
The functions are primitive recursive, including when $n$ is a parameter-dependent primitive-recursive bound. We use $\forall i<n$ and $\exists i<n$ merely as abbreviations in the target object language; their bounded-quantifier spelling can be implemented by $i<n\to\dots$ and $i<n\land\dots$.
\begin{lemma}[Finite universal reflection]\label{lem:universal}
For every fixed $f$, iPRA derives, uniformly in $n,\vec x$,
\[U_f(n,\vec x)=1\quad\longleftrightarrow\quad \forall i<n\,[f(i,\vec x)=1].\]
More precisely, the forward direction is derivable using the defining equations and bounded arithmetic, and the reverse direction is a theorem of iPRA, not an invocation of bounded induction as an axiom.
\end{lemma}
\begin{proof}
First prove $U_f(n)\le1$ by quantifier-free induction on $n$. For the forward direction, derive in the object theory the quantifier-free invariant
\[\forall i<n\,[f(i)=1] \text{ indirectly through }\ U_f(n)=1\implies f(i)=1\text{ for each parameter }i<n.\]
Formally fix $i$ and induct on $n$ on the quantifier-free formula $(i<n\wedge U_f(n)=1)\to f(i)=1$. In the successor case, split the decidable comparison $i<n$ versus $i=n$ (and $i>n$ is impossible under $i<n+1$). If $i<n$, $U_f(n+1)=U_f(n)f(n)=1$ forces both factors to be $1$, and the induction hypothesis gives $f(i)=1$. If $i=n$, use the factor $f(n)$. This is an induction on a \emph{quantifier-free} formula with a free parameter $i$. For the reverse direction, use another quantifier-free induction on $n$ with the local hypothesis $\forall i<n\ f(i)=1$ to obtain $U_f(n)=1$: in the successor step instantiate at $i=n$, restrict the premise to $i<n$, and use the induction hypothesis. Strictly speaking, the last natural-language induction carries a quantified premise; to avoid non-quantifier-free induction, instead prove, by quantifier-free induction with free parameter $n$, the more general pointwise algebraic lemma
\[\bigl(U_f(n)\neq1\bigr)\to \exists i<n\,[f(i)=0]\]
by a \emph{bounded-search witness}, encoded as the characteristic equation for a primitive-recursive search operation. Then a hypothetical $\forall i<n\,f(i)=1$ contradicts the concrete witness. Equivalently, derive this from the finite-disjunction reflection lemma below. No bounded-formula induction axiom is used.
\end{proof}

\begin{lemma}[Least failure witness]\label{lem:search}
For every fixed Boolean $f$, there are primitive-recursive terms $\mu_f(n,\vec x)$ and $h_f(n,\vec x)$ such that iPRA proves
\begin{align*}
h_f=1&\leftrightarrow \exists i<n\,[f(i)=0],\\
h_f=1&\to\mu_f<n\land f(\mu_f)=0,\\
(h_f=1\land j<\mu_f)&\to f(j)=1,\\
h_f=0&\to \mu_f=n.
\end{align*}
All parameters are implicit; $\mu_f$ is the least zero or the sentinel $n$.
\end{lemma}
\begin{proof}
Define $\mu_f(0)=0$ and $\mu_f(n+1)=\mu_f(n)$ if $\mu_f(n)<n$, $n$ if $\mu_f(n)=n$ and $f(n)=0$, and $n+1$ otherwise. Comparison and selection are primitive recursive; define $h_f(n)=\operatorname{ltb}(\mu_f(n),n)$. A joint quantifier-free induction establishes: $\mu_f(n)\le n$; whenever $\mu_f(n)<n$ its $f$-value is zero; and for each free parameter $j<\mu_f(n)$ its $f$-value is one. A second quantifier-free induction on $n$ with an additional free $i$ proves $i<n\land f(i)=0\to\mu_f(n)\le i$. The claimed existential equivalence follows by exhibiting $\mu_f$ in one direction and applying the latter bound to a given $i$ in the other. The familiar \texttt{IPRADerivations.Core.Search} already proves precisely this style of least-zero specification for its chosen convention; reuse that theorem after matching $f$ and parameter arities.
\end{proof}

\begin{proposition}[Boolean aggregation]\label{prop:agg}
In iPRA,
\[E_f(n)=1\leftrightarrow \exists i<n\,[f(i)=1],\qquad U_f(n)=1\leftrightarrow \forall i<n\,[f(i)=1].\]
\end{proposition}
\begin{proof}
The Boolean complement $1-f$ is itself primitive recursive and Boolean. A quantifier-free induction on $n$ using the recursion equations proves $E_f(n)=1-U_{1-f}(n)$; this uses distributivity and the two Boolean cases. The search lemma for $1-f$ establishes that $E_f(n)=1$ iff a witness has $f(i)=1$, and also supplies the least witness. For $U_f$, Booleanity gives $U_f=0$ or $1$. The finite search lemma for $f$ gives $U_f=0$ iff there is an index $i<n$ with $f(i)=0$: prove the equality $U_f=0$ iff $E_{1-f}=1$ by the same recursion. Universal reflection now follows by intuitionistic reasoning, using the concrete least-zero witness for the backward direction. Notice that decidability, not unrestricted excluded middle, is what makes $\neg\exists i<n\, [f(i)=0]$ equivalent to $\forall i<n\,[f(i)=1]$ and the latter to $U_f=1$.
\end{proof}

\section{The bounded reflection compiler}
\begin{definition}[Recursive Booleanization]
For each fixed $A\in\bfam$ define $e_A$ structurally: $e_{s=t}=\operatorname{eqb}(s,t)$; map connectives to \(\operatorname{and},\operatorname{or},\operatorname{not}\); set
\[e_{\forall y<B\,A}=U_{e_A}(B,\vec x),\qquad e_{\exists y<B\,A}=E_{e_A}(B,\vec x).\]
The binder is compiled via a primitive-recursive term that takes the bound variable as its first argument. The recursive construction is at the meta-level on the \emph{fixed finite syntax tree}, not a claim that a single primitive-recursive function evaluates all arbitrary encoded primitive-recursive programs.
\end{definition}
\begin{theorem}[Bounded reflection and decision]\label{thm:reflection}
For every fixed bounded formula $A(\vec x)\in\bfam$ there is a primitive-recursive function $e_A(\vec x)\in\{0,1\}$ for which iPRA derives
\[A(\vec x)\longleftrightarrow e_A(\vec x)=1,\qquad A(\vec x)\vee\neg A(\vec x).\]
The derivations are uniform in the parameters and transport under substitution of primitive-recursive terms.
\end{theorem}
\begin{proof}
Structural induction on $A$ in the metatheory constructs the term and a pair of object-theory derivations. Atomic equality is reflected by the checked arithmetic equality test. In the conjunction and disjunction cases, use the Boolean equations and the induction hypothesis. For negation, the induction hypothesis reduces $\neg A$ to $e_A=0$ because the bit is provably $0$ or $1$. At a bounded existential or universal, the recursive body is already Boolean, and Proposition~\ref{prop:agg} supplies the equivalence. Its proof uses only quantifier-free induction on auxiliary characteristic equations. Since equality with $1$ for a Boolean term is decidable, the final excluded middle is constructively derivable for $A$. Term substitution commutes with the recursively built arithmetic functions by the usual substitution lemma and instantiation of the displayed object-theory derivations. This proves an effective \emph{schema}: each finite syntax tree yields a finite proof. No single internally total interpreter for all PR programs is assumed.
\end{proof}

\begin{theorem}[Canonical guarded bounded choice]\label{thm:choice}
Let $R(\vec x,y)$ be a fixed member of $\bfam$, and $B(\vec x)$ a primitive-recursive bound. There are primitive-recursive functions $w(\vec x)$ and $h(\vec x)\in\{0,1\}$ such that iPRA proves
\begin{align*}
h=1&\leftrightarrow\exists y<B\ R(\vec x,y),\\
h=1&\to w<B\land R(\vec x,w),\\
h=0&\to w=B,\\
(h=1\land z<w)&\to\neg R(\vec x,z).
\end{align*}
Consequently $\forall \vec x\,\exists y<B(\vec x)\,R(\vec x,y)$ entails $\forall \vec x\,R(\vec x,w(\vec x))$ in iPRA, with a concrete PR witness.
\end{theorem}
\begin{proof}
Compile $R$ to $e_R$ by Theorem~\ref{thm:reflection}. Apply least-zero search to $f=1-e_R$ with bound $B$. Its sentinel and leastness laws yield all four statements. Under the universal existence premise, the characteristic success flag equals one at each $\vec x$; eliminate the guard and use universal introduction. This is not unrestricted choice: the \emph{explicit primitive-recursive bound} is an essential hypothesis.
\end{proof}

\section{Parameter-dependent finite covers and local-to-global rules}
\begin{theorem}[Guarded finite cover]\label{thm:cover}
Suppose $C(i,x)$ is bounded, $N(x)$ is primitive recursive, and $P(x)$ is an arbitrary formula (not assumed decidable). If iPRA proves
\[\forall x\,\exists i<N(x)\,C(i,x)\qquad\text{and}\qquad\forall x\forall i<N(x)\,[C(i,x)\to P(x)],\]
then iPRA proves $\forall x\,P(x)$. Moreover, the proof may be specialized to the canonical primitive-recursive witness $i=w(x)$, obtained from Theorem~\ref{thm:choice}.
\end{theorem}
\begin{proof}
Fix $x$. The first premise implies $h(x)=1$ by bounded reflection; guarded choice yields $w(x)<N(x)$ and $C(w(x),x)$. Instantiate the second premise at $w(x)$, and apply modus ponens. Finally generalize $x$. The inference uses neither a decision procedure for $P$ nor a proof of $P$ by induction on a quantified invariant. The canonical-witness refinement makes the proof stable under further object-language term substitution.
\end{proof}

\begin{theorem}[Composing two bounded searches]\label{thm:composition}
Suppose $R(x,y)$ and $S(x,y,z)$ are bounded and $B(x),D(x,y)$ are primitive recursive. If iPRA proves
\[\forall x\,\exists y<B(x)\ R(x,y),\quad \forall x\forall y<B(x)\,[R(x,y)\to\exists z<D(x,y)\ S(x,y,z)],\]
then there exist primitive-recursive witness functions $w(x),v(x)$ such that iPRA proves
\[\forall x\,[R(x,w(x))\land S(x,w(x),v(x))].\]
\end{theorem}
\begin{proof}
Apply guarded choice to $R$ to get $w(x)$ and proof of $R(x,w(x))$. Specialize the second premise at $y=w(x)$ and use the provable inequality $w(x)<B(x)$. Apply guarded choice for $S$ at parameters $(x,w(x))$. Composition of primitive-recursive functions produces $v(x)$, and conjunction introduction completes the proof. Repeating this argument for any fixed finite number of stages produces a finite, statically nested witness pipeline; no countable dependent choice is used.
\end{proof}

\section{A worked instance: bounded polynomial certificates}
Let $p(x,y)$ be a fixed integer polynomial, interpreted as equality $p_+(x,y)=p_-(x,y)$ of natural-number polynomials. Fix a PR bound $B(x)$ and write
\[R(x,y):\equiv p_+(x,y)=p_-(x,y),\qquad\mathrm{sol}(x)=\mu_{y<B(x)} [p_+(x,y)=p_-(x,y)].\]
The function $\mathrm{sol}$ is primitive recursive with the sentinel $B(x)$. Theorem~\ref{thm:choice} gives in iPRA
\[\exists y<B(x)\,p(x,y)=0\ \leftrightarrow\ \mathrm{sol}(x)<B(x),\]
\[\mathrm{sol}(x)<B(x)\ \to\ p(x,\mathrm{sol}(x))=0\ \wedge\ \forall j<\mathrm{sol}(x)\,p(x,j)\ne0.\]
Thus bounded Diophantine witness extraction is unconditional, for a specified polynomial and bound, even though no analogous uniform unbounded Diophantine solver exists. One can compose this with an explicitly bounded verification of polynomial identities, resultant certificates, or finite matrices. This example exposes the exact interface needed from the repository's \texttt{Core.Search}: substitute the candidate equality bit as the zero-test body and instantiate its \texttt{search\_found} and \texttt{search\_minimal} theorems.

\section{Where the method stops}
\begin{warning}[Do not confuse three layers]
(1) A metatheoretic Lean theorem about a function on standard $\mathbb N$ is not itself an iPRA derivation. (2) A proof packet containing a Lean \texttt{sorry} or a comment is not a kernel-checked derivation. (3) A Booleanizer for each \emph{fixed bounded formula} does not imply a single primitive-recursive universal evaluator for arbitrary primitive-recursive codes.
\end{warning}
\begin{warning}[Absence of effective bounds]
From an unbounded $\forall x\exists y\,R(x,y)$, one cannot simply infer a primitive-recursive bound or witness. In particular the universal bounded-choice theorem gives no general algorithm for Hilbert's tenth problem. A classical compactness argument establishing some infinite object need not furnish a uniform finite PR modulus, and a finite family of approximations $A(n)$ for every $n$ does not establish an unqualified infinite limit assertion unless a suitable coding, quantitative modulus, and target formula are provided.
\end{warning}
\begin{warning}[Proof-theoretic scope]
The present arguments use the defining equations and quantifier-free induction in the \emph{ordinary} iPRA language. The \texttt{Bridge} currently exposes closed-sentence translation; lifting the entire framework to Foundation arbitrary open sequents needs the still-unfinished open-context and round-trip work described in \texttt{UNIFICATION-PLAN.md}. Likewise, formalizing arbitrary bounded formula syntax and its compiler is a finite metatheoretic project; each compiled result must end with an explicit \texttt{Prf PRAAx} object to count as an iPRA theorem.
\end{warning}

\section{Lean formalization architecture}
We propose three independent modules. \texttt{BoundedReflection/Semantics.lean} defines a lightweight mathematical evaluator and proves correctness for each finite bounded syntax tree. This is a \emph{testing layer}, not the final evidence. \texttt{BoundedReflection/Object.lean} builds iPRA terms and \texttt{Prf PRAAx} derivations from fixed schema constructors. \texttt{BoundedReflection/Applications.lean} derives guarded choice, cover and search composition for concrete arities and reuses the checked \texttt{IPRADerivations.Core.Search}.

For a rapid integration, first construct only the \emph{fixed-formula API}: given $f:\texttt{PTm 2}$ and a proven equation \texttt{f(i,x)=0} encoding a predicate, derive the search theorem at variable bound $N$ by invoking \texttt{Search.search\_found} and \texttt{Search.search\_minimal}. Then use the checked finite sums / Boolean operations to implement a finite universal aggregator. Finally introduce a typed inductive \texttt{BForm n} of formulas with terms \texttt{PTm n} and bounded binders, and define mutually structural compilation functions \texttt{compileTerm} and \texttt{compileProof}. In the quantifier cases, use an explicit arity-indexed abstraction of the body into a PR symbol; avoid treating an arbitrary Lean closure as a PR code. The most important code obligation is the substitution/binder lemma with variable shifts; write this before automating proofs.

The supplied \texttt{LeanOutline.lean} is an intentionally uncompiled specification file containing the mathematical semantic prototype and labeled proof obligations. It is \emph{not} presented as checked object-code. The companion \texttt{finite\_reflection.py} is executable test code for finite bounds, not a proof checker or a trusted Lean artifact.

\section{Recommended integration and verification milestones}
\begin{enumerate}[leftmargin=*]
\item Reuse the current \texttt{IPRADerivations.Core.Search} and \texttt{Core.Order} instead of defining a competing least-witness theory. Compare conventions and document the arities and sentinel.
\item Prove an open object-theory test instance for $R(x,y)\equiv y+y=x$ with a variable bound $B(x)=x+1$ and a concrete \texttt{Prf PRAAx} conclusion; show the negative branch with sentinel.
\item Add the universal-conjunction PR symbol and its equation, then prove \texttt{U\_reflect} under arbitrary hypotheses $\Gamma$ by object-theory quantifier-free induction.
\item Construct the inductive syntax compiler, prove substitution and two object-derivation directions for every constructor. Check with \texttt{lake build IPRADerivations} and a standalone focused module.
\item Add regression tests: zero bound, zero-free search, nested alternating quantifiers, variable parameter bounds, and a two-stage witness pipeline. Audit \texttt{\#print axioms} of each exported theorem. Keep the Lean compiler's metatheoretic semantic tests separate from the object-theory proof certificate.
\item Only after the existing bridge gains open-context transport should the result be exported as a general Foundation open-sequent theorem; meanwhile closed universal closures can use the checked closed bridge where applicable.
\end{enumerate}

\section{Comparison with the existing framework packets}
The repository already contains numerous finite-query, finite-use, staged-contract, categorical, and observation frameworks under \texttt{docs/proofs/foundations/}; some propose essentially similar reusable combinators. The differentiator here is a deliberately small \emph{theorem-level contract}: a fixed bounded syntax class; explicit PR upper bounds; a Boolean evaluator accompanied by two object-theory derivations; a canonical sentinel witness; and a precise statement of what is \emph{not} compiled. It does not claim priority or a novel proof of bounded decidability. Its intended contribution is to make these standard principles a sound, reusable, auditable interface to existing checked arithmetic. A production implementation should first deduplicate it against the above framework packets and the checked \texttt{Core.Search}.

\begin{thebibliography}{9}
\bibitem{repo} V. Reshetnikov, \emph{ProveIt}, \url{https://github.com/VladimirReshetnikov/ProveIt/tree/main/Logic/PrimitiveRecursiveArithmetic}, accessed 16 Nov 2026.
\bibitem{plan} \emph{iPRA unification plan}, \url{https://github.com/VladimirReshetnikov/ProveIt/blob/main/Logic/PrimitiveRecursiveArithmetic/UNIFICATION-PLAN.md}.
\bibitem{deriv} \emph{IPRADerivations README}, \url{https://github.com/VladimirReshetnikov/ProveIt/blob/main/Logic/PrimitiveRecursiveArithmetic/Derivations/README.md}.
\bibitem{search} \emph{Core/Search.lean}, \url{https://github.com/VladimirReshetnikov/ProveIt/blob/main/Logic/PrimitiveRecursiveArithmetic/Derivations/Lean/IPRADerivations/Core/Search.lean}.
\bibitem{bridge} \emph{IPRABridge README}, \url{https://github.com/VladimirReshetnikov/ProveIt/blob/main/Logic/PrimitiveRecursiveArithmetic/Bridge/README.md}.
\end{thebibliography}
\end{document}
